\subsection{Gestione conversazioni}

Nella parte sinistra della schermata trovi l'area delle conversazioni, che ti permette di organizzare e gestire tutte le interazioni con la chat.

% TODO: inserire figura area conversazioni nella schermata principale

Ogni conversazione rappresenta uno scambio separato, cos\`i puoi lavorare su pi\`u richieste senza confonderti.

\subsubsection{Creazione nuova conversazione}

Per iniziare una nuova conversazione puoi:
\begin{itemize}[itemsep=5pt, parsep=5pt, label=$\scriptstyle\bullet$]
	\item Cliccare sul pulsante \vr{+} in alto
	\item Selezionare la voce \vr{Nuova conversazione} nell'elenco
\end{itemize}

% TODO: inserire figura creazione nuova conversazione

In entrambi i casi verr\`a creata una nuova chat vuota, pronta per inserire il primo messaggio.

Accanto al pulsante di creazione trovi anche un'icona che ti permette di aprire o chiudere il pannello delle conversazioni, utile per avere pi\`u spazio nella schermata quando ne hai bisogno.

% TODO: inserire figura icona apertura e chiusura pannello conversazioni

\subsubsection{Selezione conversazione}

Nell'elenco vengono mostrate tutte le conversazioni disponibili.\\
Per aprirne una, ti basta cliccarci sopra: la conversazione selezionata diventa attiva e puoi vedere tutti i messaggi nella chat centrale.

% TODO: inserire figura selezione conversazione

La conversazione attiva \`e evidenziata, cos\`i puoi capire subito su quale stai lavorando.\\
Puoi passare da una conversazione all'altra in qualsiasi momento, senza perdere lo storico dei messaggi.

\subsubsection{Gestione conversazione}

Ogni conversazione pu\`o essere gestita tramite il menu accessibile con l'icona a tre puntini, accanto al suo nome.

% TODO: inserire figura menu gestione conversazione

Da qui puoi:
\begin{itemize}[itemsep=5pt, parsep=5pt, label=$\scriptstyle\bullet$]
	\item Rinominare la conversazione, utile per darle un nome pi\`u chiaro e riconoscerla facilmente
	\item Eliminare la conversazione, se non ti serve pi\`u
\end{itemize}

Prima di eliminare una conversazione, il sistema ti chieder\`a una conferma, perch\'e tutti i messaggi verranno cancellati definitivamente.

% TODO: inserire figura conferma eliminazione conversazione

Puoi quindi decidere se:
\begin{itemize}[itemsep=5pt, parsep=5pt, label=$\scriptstyle\bullet$]
	\item Confermare, eliminando la conversazione
	\item Annullare, mantenendo tutto com'\`e
\end{itemize}
